% (C) 2026 Ralf Stephan, in collaboration with Claude Code.
% Released under CC0 1.0 Universal (public-domain dedication).
% See https://creativecommons.org/publicdomain/zero/1.0/

\documentclass[11pt]{article}
\usepackage[T1]{fontenc}
\usepackage[utf8]{inputenc}
\usepackage{lmodern}
\usepackage[margin=2.5cm]{geometry}
\usepackage{amsmath,amssymb}
\usepackage{braket}
\usepackage[hidelinks]{hyperref}

% Lean names and paths: verbatim, breakable at . _ /
\DeclareUrlCommand\lean{\urlstyle{tt}}
\def\UrlBreaks{\do\.\do\_\do\/}

\title{Open problems:\\ Equivalence of sup-norm, $L^1$, and $L^2$ seminorms on the Schwartz
  space}
\author{}
\date{}

\begin{document}
\maketitle

\noindent
Open problems raised by the technical brief in this directory, and by its Lean formalization
(see \lean{formalization-log.tex}).

\section*{From the brief}
\begin{enumerate}
\item \textbf{Finite moments imply Schwartz.} ``Finiteness of all such expectation values is
  equivalent to membership in $\mathcal{S}(\mathbb{R})$'', and $\mathcal{S}(\mathbb{R})$ ``is
  exactly the space of states for which all polynomial moments of position and momentum are
  well-defined''. Only the bound (10) is formalized. The converse needs a precise statement:
  which orderings of $X^\alpha$ and $P^\beta$ count, and what ``well-defined'' means for $f$
  outside their domains, since $(\|f\|^{L^2}_{\alpha,\beta})^2 =
  \braket{X^\alpha\partial^\beta f|X^\alpha\partial^\beta f}$ is not of the form
  $\braket{f|X^\alpha P^\beta|f}$. (Brief, Abstract, Remark after eq.~(10) and Sec.~4; log,
  ``Differences from the brief'': not formalized.)
\item \textbf{General $L^p$ seminorms.} ``The same arguments extend to $L^p$ seminorms for any
  $1 \le p \le \infty$'', and all $L^p$-based Schwartz seminorm families generate the same
  topology. Not formalized. (Brief, Remark ``General $L^p$''; log, ``Differences from the
  brief'': not formalized.)
\end{enumerate}

\section*{From the formalization}
\begin{enumerate}
\item None found.
\end{enumerate}

\end{document}
